\documentclass[11pt]{article}
\usepackage[utf8]{inputenc}
\usepackage[T1]{fontenc}
\usepackage{lmodern}
\usepackage[margin=1in]{geometry}
\usepackage{amsmath,amssymb,amsthm,mathtools}
\usepackage{booktabs}
\usepackage{enumitem}
\usepackage[colorlinks=true,linkcolor=black,citecolor=black,urlcolor=black]{hyperref}
\usepackage{longtable}
\usepackage{array}
\usepackage{graphicx}
\usepackage{microtype}
\usepackage[most]{tcolorbox}
\tcbuselibrary{skins,breakable}
\usepackage{setspace}
\setstretch{1.05}
\setlength{\parskip}{0.65em}
\setlength{\parindent}{0pt}

\newtheorem{theorem}{Theorem}[section]
\newtheorem{definition}[theorem]{Definition}
\newtheorem{remark}[theorem]{Remark}
\newtheorem{proposition}[theorem]{Proposition}
\newtheorem{lemma}[theorem]{Lemma}

\newcommand{\coderef}[2]{\texttt{#1}\,(\texttt{#2})}

\newtcolorbox{apfbox}{
  breakable, enhanced,
  colback=black!3, colframe=black!30,
  boxrule=0.6pt, arc=2pt,
  left=8pt, right=8pt, top=8pt, bottom=8pt
}

\newtcolorbox{wayfinder}{
  breakable, enhanced,
  colback=blue!3, colframe=blue!40,
  boxrule=0.7pt, arc=2pt,
  left=8pt, right=8pt, top=8pt, bottom=8pt,
  title={\textbf{Wayfinder}},
  fonttitle=\bfseries,
  coltitle=blue!50!black, colbacktitle=blue!10
}

\newtcolorbox{bridgebox}{
  breakable, enhanced,
  colback=orange!3, colframe=orange!50,
  boxrule=0.7pt, arc=2pt,
  left=8pt, right=8pt, top=8pt, bottom=8pt,
  title={\textbf{Bridge}},
  fonttitle=\bfseries,
  coltitle=orange!60!black, colbacktitle=orange!10
}

\newtcolorbox{statusbox}{
  breakable, enhanced,
  colback=red!3, colframe=red!40,
  boxrule=0.7pt, arc=2pt,
  left=8pt, right=8pt, top=8pt, bottom=8pt,
  title={\textbf{Status fence}},
  fonttitle=\bfseries,
  coltitle=red!60!black, colbacktitle=red!10
}

\title{Paper 6 --- REPRESENTATIONS (Dynamics \& Geometry):\\
Dynamics and Geometry as Least-Cost Admissible Reallocation\\[2pt]
\large Variational Form, Metric Structure, Einstein Equations, and Their Regime Boundaries}
\author{E. Brooke}
\date{April 2026 (v2.0-PLEC)}

\begin{document}
\maketitle

\begin{abstract}
This paper derives dynamics and geometry as representational structures that emerge
when admissibility constraints support smooth, locally additive bookkeeping over a
connected admissible path class. Equations of motion arise only in restricted
regimes; outside these regimes --- at saturation, branching, record locking, or
horizons --- variational dynamics ceases to be meaningful, though admissibility
accounting remains well-defined. Geometry emerges by the same route: distance is
minimal correlation cost, curvature arises from capacity gradients, and Einstein's
equations appear as the unique closure relation in four dimensions (Lovelock
uniqueness). The v2.0-PLEC release reframes the derivation around the Principle
of Least Enforcement Cost (PLEC): admissibility fixes the licensable continuation
class; PLEC selects the realized continuation as a least-cost path within that
class; failure of dynamical or geometric representation is recast as one of five
formally distinct regime-exit types. Imported closure theorems (HKM, Malament,
Lovelock) are tagged as imports closing an APF-prepared admissible class. The
dark-matter-as-$C_{\mathrm{ext}}$ identification is fenced as a structural
proposal with empirical consequences, not as a theorem-level closure.
\end{abstract}

\begin{wayfinder}
\textbf{Operational primitives used in this paper.}
\begin{itemize}[leftmargin=1.5em,nosep]
    \item \textbf{Admissible class $\mathcal{A}_{\Gamma}$:} continuations
    physically licensable in context $\Gamma$ under A1.
    \item \textbf{Enforcement-cost functional $E_{\Gamma}$:} total capacity burden
    of a continuation; on paths, an action functional $K[q] = \int L\,dt$.
    \item \textbf{PLEC selection $q^{\ast}$:} realized continuation,
    $q^{\ast} \in \arg\min_{q \in \mathcal{P}_{\Gamma}} K[q]$.
    \item \textbf{Regime R:} smooth, locally additive enforcement cost on a
    connected admissible path space without saturation. Variational and geometric
    representation valid here.
    \item \textbf{Regime exit:} failure of one or more Regime R conditions.
    Catalogued in five formally distinct types (\S\ref{sec:exit-taxonomy}).
\end{itemize}
This paper writes Euler--Lagrange dynamics and Lovelock geometry as the realized
PLEC selectors on the Regime R admissible class, and writes saturation, branching,
record locking, horizons, singularities, and topology change as different formal
exits from that representational class.
\end{wayfinder}

\section*{Shared Canonical Frame (Papers 5--6)}

\begin{apfbox}
\textbf{Meaning, Admissibility, Realization, Regime Exit.}
\begin{enumerate}[leftmargin=1.6em]
    \item \textbf{Meaning.} A distinction is physically meaningful iff it is
    robustly enforceable under admissible perturbations.
    \item \textbf{Admissibility.} A configuration, path, or representation is
    admissible iff its total enforcement demand does not exceed available capacity.
    \item \textbf{Realization.} When the admissible class is nonempty and the
    enforcement-cost functional is well-defined on it, realized physical structure
    is selected as the least-enforcement-cost member of that admissible class.
    \item \textbf{Regime exit.} When the conditions required for comparison by
    smooth finite cost fail, the representation exits its domain, while the deeper
    admissibility constraints remain in force.
\end{enumerate}
\end{apfbox}

\section*{Part I: Dynamics}

\section{Introduction}

\subsection{The Standard View}
In conventional physics, dynamics and geometry are taken as primitive. Equations
of motion are postulated; spacetime is assumed as an arena. The questions ``why
do physical laws take variational form?'' or ``why is gravity geometric?'' are
rarely asked at the level of structure.

\subsection{The Constraint-First Inversion}
This paper asks: under what conditions do equations of motion exist at all? Under
what conditions does geometric description become meaningful? What happens when
those conditions fail? The central claim is that dynamics and geometry are
\emph{regime-dependent representations of admissible correlation reallocation}.
They emerge as PLEC selectors when correlation capacity admits smooth
redistribution, enforcement cost is locally additive, and saturation boundaries
are not encountered. They exit cleanly in five distinct ways when those
conditions fail.

\section{Canonical PLEC Schema (Path Form)}\label{sec:plec}

\begin{definition}[Admissible path class]
Let $\Gamma$ denote a physical context, interface, or regime. Let $\mathcal{P}_{\Gamma}$
denote the set of admissible paths joining fixed endpoints in configuration space
in context $\Gamma$.
\end{definition}

\begin{definition}[Enforcement-cost functional on paths]
Let $K_{\Gamma} : \mathcal{P}_{\Gamma} \to [0,\infty)$ assign the total accumulated
enforcement cost along a path $q(t)$ in $\Gamma$. In Regime~R (\S\ref{sec:regimeR}),
$K_{\Gamma}[q] = \int L(q,\dot q,t)\,dt$ for a local cost density $L$.
\end{definition}

\begin{definition}[PLEC selection on paths]
A realized continuation $q^{\ast}$ in context $\Gamma$ is selected by
\[
q^{\ast} \in \operatorname*{arg\,min}_{q \in \mathcal{P}_{\Gamma}} K_{\Gamma}[q].
\]
When the minimizer is unique up to the relevant equivalence relation, the realized
continuation is unique up to that equivalence.
\end{definition}

\section{Admissibility Does Not Yet Determine Realized Evolution}

Admissibility determines which correlation structures are physically licensable.
It specifies the class of jointly enforceable configurations and rules out those
that overspend finite capacity. In this sense, admissibility provides necessary
conditions for physical meaningfulness. But admissibility alone does not yet
determine realized evolution. Given an admissible initial state, there may exist
multiple distinct admissible continuations. The admissibility framework therefore
fixes the \emph{allowed class} of continuations without, by itself, selecting a
unique realized one.

This distinction is fundamental. The framework must separate three questions
that are often run together:
\begin{enumerate}[leftmargin=1.6em]
    \item \textbf{Meaning:} when is a distinction physically meaningful at all?
    \item \textbf{Admissibility:} which continuations are jointly enforceable
    under finite capacity?
    \item \textbf{Realization:} which admissible continuation is physically
    realized?
\end{enumerate}
Papers~1--4 answer the first two questions at the level of enforceable
distinction and capacity-bounded admissibility. The task of the present paper is
to identify the additional regime conditions under which the third question
becomes mathematically well-posed. \emph{Dynamics is not automatic from
admissibility; it arises only when the admissible class carries enough structure
to support a least-cost selection principle.}

\section{Regime R --- When Least-Cost Selection Becomes Meaningful}\label{sec:regimeR}

Let $\mathcal{A}_{\Gamma}$ denote the class of admissible continuations in
context $\Gamma$. In general, $\mathcal{A}_{\Gamma}$ may be disconnected,
non-smooth, or terminate at saturation boundaries. In such cases, there need be
no meaningful notion of a realized path selected from within
$\mathcal{A}_{\Gamma}$. A dynamical representation becomes available only in a
restricted regime:

\begin{enumerate}[label=(R\arabic*),leftmargin=1.8em]
    \item Enforcement cost varies smoothly over admissible correlation sets.
    \item Cost is locally additive over interfaces.
    \item Admissible continuations form a connected path space.
    \item No saturation boundary is encountered along the path.
\end{enumerate}

When all four conditions hold, the admissible class supports a well-defined cost
functional on paths. In that regime, physical realization can be expressed as
\emph{least-cost admissible reallocation}: among admissible continuations, the
realized path is the one minimizing accumulated enforcement cost.

This is the variational content of the framework. It is not an additional
dynamical postulate imposed from outside. It is the structural consequence of
asking for realized evolution only where the admissible class admits comparison
by smooth finite cost.

\section{Least-Cost Admissible Paths and Euler--Lagrange Form}\label{sec:EL}

In Regime~R, the accumulated enforcement cost along a path defines an action
functional
\[
K[q] = \int L(q,\dot q,t)\,dt,
\]
where $L$ is the local enforcement-cost density. The physical claim is then:

\begin{apfbox}
\textbf{Selection Principle (PLEC in Regime R).} The realized continuation is a
least-enforcement-cost member of the admissible path class:
\[
q^{\ast} \in \operatorname*{arg\,min}_{q \in \mathcal{P}_{\Gamma}} K[q].
\]
\end{apfbox}

Once this selection principle is available, standard variational reasoning
applies. Extremal admissible paths satisfy the Euler--Lagrange equations
\[
\frac{d}{dt}\left(\frac{\partial L}{\partial \dot q}\right) - \frac{\partial L}{\partial q} = 0.
\]
These equations are therefore not primitive laws. They are the coordinate
expression of least-cost selection within the admissible class.

The usual conservation laws follow in the same way. Whenever the cost functional
has a continuous symmetry, Noether's theorem identifies the corresponding
conserved quantity. Energy conservation expresses invariance of enforcement cost
under time translation; momentum conservation expresses invariance under spatial
translation. The standard dynamical formalism appears as a \emph{representation of
admissible least-cost realization}, not as an independent ontological layer.

\section{Where Least-Cost Dynamical Representation Fails}\label{sec:dyn-fails}

These are not breakdowns of admissibility itself. They are breakdowns of the
hypotheses under which realized evolution can be represented as least-cost
variation over a smooth admissible path class. When Regime~R fails,
admissibility accounting remains valid, but Euler--Lagrange form ceases to be
meaningful.

\begin{center}
\begin{tabular}{@{}lll@{}}
\toprule
Mode & Failed hypothesis & Consequence \\
\midrule
Saturation & Collapse of admissible variation & No nontrivial local variation remains \\
Branching & Minimizer nonuniqueness & No single effective equation of motion \\
Record locking & Change of admissible class & Exit to record-dominated description \\
Horizons & Failure of admissible extendibility & Dynamical inaccessibility from within regime \\
\bottomrule
\end{tabular}
\end{center}

The important point is that the representational layer fails in different ways
depending on which hypothesis breaks. There is no single ``failure of dynamics'';
there are several distinct failures of least-cost selection structure, each with
its own structural fingerprint.

\textbf{On branching.} Branching is not merely the existence of multiple
admissible continuations. It is the failure of uniqueness for the least-cost
selector on the admissible class. When multiple inequivalent minimizers exist,
admissibility remains intact, but the representational layer no longer supports
a unique effective dynamics. The correct conclusion is not that physical law
fails, but that the given regime no longer compresses realized evolution into a
single variational trajectory.

\section{Taxonomy of Regime Exit}\label{sec:exit-taxonomy}

To unify the representation failures discussed across Papers~5 and~6, we
distinguish five structurally different kinds of regime exit. This taxonomy
provides a common project-level vocabulary.

\begin{figure}[h]
\centering
\includegraphics[width=0.82\textwidth]{Fig_Admissible_States_Diagram.pdf}
\caption{\textbf{Admissible states and correlation-capacity space.} The upper
bulb is the admissible class $\mathcal{A}_{\Gamma}$ (Paper~5/6 \S2); its
boundary is where admissibility fails (flanking shaded regions, labeled
``Not Admissible''). The downward arrow is Regime~R (\S\ref{sec:regimeR}):
time emerges as the direction of cost-functional integration along an
admissible path, not as a primitive coordinate. At the narrowing near the
bottom, the locking interface marks the failure mode of unique PLEC
selection. Branching into three tubes visualizes Type~II
(\coderef{check\_Regime\_exit\_Type\_II}{plec.py}) — non-unique minimizers
of the cost functional. The step to three distinct downstream classes (the
``exclusive correlation histories,'' each associated with its own
record-locked sector) visualizes Type~III
(\coderef{check\_Regime\_exit\_Type\_III}{plec.py}) — change of admissible
class at the locking interface, the mechanism identified in
\coderef{check\_L\_irr}{core.py} whereby correlation pathways become
inaccessible to local operations (not infinite-cost; locally
unrecoverable while remaining globally admissible). Probability assignments
are defined only on the entropy slice downstream of locking
(\coderef{check\_L\_Gleason\_finite}{supplements.py}); before branching,
the system is in the pre-statistical coherent regime of Paper~5.
Types~I (saturation collapse), IV (loss of smooth structure), and V (pure
representational redundancy) are not branching-type exits and are not
depicted; they are covered separately in \S\ref{sec:dyn-fails} and
\S\ref{sec:geo-fails}. The volume/area imagery at the bottom
(``capacity is a volume, entropy is an area'') is a visual anchor for
the Bekenstein area-law relation \coderef{check\_T\_Bek}{gravity.py}, not
a literal 2--3 dimensional assertion: the admissible class is generally
higher-dimensional and the entropy slice is an interface whose
area-scaling is the physical content.}
\label{fig:admissible-states}
\end{figure}

Figure~\ref{fig:admissible-states} visualizes the full arc of this section
at a glance. A reader tracking the five types below can consult the figure
as the running example for Types~II and~III; the remaining types (I, IV,
V) are qualitatively different and need no figure.

\subsection{Type I: Collapse of admissible variation}
A representational regime fails when the admissible neighborhood around a state
or path collapses, so that no nontrivial admissible variation remains. This is
the canonical saturation case. In Paper~6 language, the variational
representation fails because the local path class becomes trivial; in Paper~5
language, coherent refinement fails because no reversible alternatives remain.

\subsection{Type II: Minimizer nonuniqueness}
A regime can also fail to determine a single realized continuation even when
admissible paths still exist. This is the branching case: the admissible class
remains nonempty, but the least-cost selector fails to be unique up to the
relevant equivalence relation. The correct interpretation is not that physical
law has failed, but that the representational layer no longer supports a unique
effective dynamics.

\subsection{Type III: Change of admissible class}
Some regime exits occur because the relevant admissible class itself changes.
Measurement is the prototype. Before record formation, the relevant class is the
coherent bookkeeping class; after irreversible logging, the relevant class is the
record-locked class on $M_{\mathrm{sys}} \otimes Z_R$. The transition is
therefore not a breakdown internal to a single representational scheme, but a
transfer to a different admissible domain.

\subsection{Type IV: Loss of smooth or local structure}
A regime may remain admissible while losing the regularity assumptions required
for a geometric or variational description. Singularities, Planck-scale
discreteness, and topology change belong here. The issue is not loss of
admissibility, but loss of the smoothness, local additivity, tangent-space
structure, or chartability required for the relevant representation.

\subsection{Type V: Pure representational redundancy}
Finally, some apparent failures are not physical exits at all, but breakdowns of
a chosen representation. Gauge redundancy and nonunique coordinate descriptions
fall in this class. The admissible structure and realized minimizer remain
intact; only the descriptive coding is nonunique. This category prevents
conflating genuine regime exit with bookkeeping arbitrariness.

\subsection{Canonical use across the APF series}
\begin{center}
\begin{tabular}{@{}lll@{}}
\toprule
Phenomenon & Type & Coderef (v6.9) \\
\midrule
Paper~5 measurement & III (with I in fully locked limit) & \coderef{check\_Regime\_exit\_Type\_III}{plec.py} \\
Paper~6 saturation & I & \coderef{check\_Regime\_exit\_Type\_I}{plec.py} \\
Paper~6 branching & II & \coderef{check\_Regime\_exit\_Type\_II}{plec.py} \\
Paper~6 horizons & II / III & \coderef{check\_Regime\_exit\_Type\_II}{plec.py} / \coderef{check\_Regime\_exit\_Type\_III}{plec.py} \\
Paper~6 singularities, Planck discreteness, topology change & IV & \coderef{check\_Regime\_exit\_Type\_IV}{plec.py} \\
Gauge / coordinate ambiguity & V & \coderef{check\_Regime\_exit\_Type\_V}{plec.py} \\
\bottomrule
\end{tabular}
\end{center}
The strategic payoff is conceptual hygiene: measurement, horizons, saturation,
and geometry failure are different species of representational regime exit, all
subordinate to the deeper admissibility layer rather than independent mysteries.
As of v6.9 (2026-04-18), each type is a registered bank check with an
executable witness; the taxonomy is no longer prose but formalized
machinery.

\section*{Part II: Geometry}

\section{Correlation Cost and Distance}

\subsection{Definition of Distance}
\begin{definition}[Correlation distance]
The distance between regions $A$ and $B$ is the minimal irreducible correlation
cost required to jointly enforce coordination between them:
\[
d(A,B) \equiv \inf_{\Gamma : A \to B} \int C\,d\Gamma,
\]
where $C$ is the correlation cost density along path $\Gamma$.
\end{definition}
Correlation distance satisfies non-negativity, identity of indiscernibles, and
the triangle inequality. These are forced by enforcement-cost structure, not
assumed.

\subsection{Geodesics as Least-Cost Paths}
Paths minimizing correlation cost define geodesics. Along such paths, relational
structure is maintained with minimal capacity expenditure. Free motion is
least-cost coordination. No force law is required.

\section{The Metric Tensor from Non-Factorization (T7B)}

\subsection{Non-Factorization Implies Geometric Response}
\begin{theorem}[T7B, Metric tensor from quadratic non-factorizing response]
At shared (non-factorizing) interfaces, the external feasibility functional is
quadratic in displacement. By the polarization identity, a functional that is
(a) non-degenerate (A1: finite capacity implies no zero-cost directions),
(b) continuous (Lipschitz, from finite-bounded cost), and (c) defined on
tangent vectors at each point (from $T_{\mathrm{continuum}}$ smooth structure)
is a metric tensor $g_{\mu\nu}$.
\end{theorem}
Formal verification: \coderef{check\_T7B}{gravity.py}.

This is a definition-level identification: a non-degenerate symmetric bilinear
form on the tangent space at each point \emph{is} a metric tensor. The framework
does not postulate geometry; it computes it from quadratic cost at shared
interfaces. Tag: [P\_structural].

\subsection{Lorentzian Signature $(-,+,+,+)$}
A4 (irreversibility) implies a strict partial order on admissible
configurations. By the Hawking--King--McCarthy (HKM, 1976) theorem, a strict
partial order on events satisfying appropriate density conditions determines the
conformal structure. Combined with Malament's theorem (1977) and volume
normalization ($\Omega = 1$ by Radon--Nikodym uniqueness), this fixes the
Lorentzian signature with exactly one timelike direction.

\begin{bridgebox}
HKM and Malament are imported as closure theorems on an APF-prepared admissible
class: APF supplies the strict partial order from A4 and the volume
normalization; the imports close the conformal structure and signature
identification. APF-internal preparation checks:
\coderef{check\_L\_HKM\_causal\_geometry}{extensions.py} and
\coderef{check\_L\_Malament\_uniqueness}{extensions.py}.
Tag: [P\_structural] + [I: HKM 1976, Malament 1977].
\end{bridgebox}

\section{Spacetime Dimension $d = 4$ (T8)}

\begin{theorem}[T8, $d = 4$ as the minimal admissible spacetime dimension]
$d \le 3$ is hard-excluded: $d = 3$ has zero propagating graviton degrees of
freedom (no gravitational records, violating A4); $d = 2, 3$ lack knots for
topological stability of bound states. $d = 4$ is the minimum dimension with
propagating gravity (2 DOF), knot stability, and capacity efficiency
($C_{\mathrm{ext}}/C_{\mathrm{int}} = 0.5$). $d \ge 5$ is disfavored by A5
(genericity: extra graviton modes provide no admissibility benefit).
\end{theorem}

The internal sector uses $C_{\mathrm{int}} = 12$ (dim $SU(3)\!\times\!SU(2)\!\times\!U(1)$),
leaving $C_{\mathrm{ext}}$ for geometric degrees of freedom. In $d=4$, the
graviton has 2 propagating DOF, the minimum for A4-compatible gravitational
records. Tag: [P\_structural]. Formal verification: \coderef{check\_T8}{spacetime.py}.

\section{Curvature from Capacity Gradients}

Correlation capacity is generally non-uniform. Capacity gradients cause geodesic
deviation: paths converge in low-capacity regions (bottlenecks) and diverge in
high-capacity regions. \emph{Curvature is the geometric signature of correlation
load} --- not a postulated feature of spacetime but a derived consequence of
capacity distribution.

The stiffness parameter $G$ (Newton's constant) sets the scale of geometric
response to correlation load: (curvature change) $\sim G \times$ (load change).
$G$ is not a coupling strength but the rate at which geometry bends under
correlation burden. Structurally, $\kappa \sim 1/C_{\ast}$ where $C_{\ast}$ is
the total geometric capacity (T10). Tag: [P\_structural], scaling only.
Formal verification: \coderef{check\_T10}{gravity.py}.

\section{Einstein's Equations from $\Delta_{\mathrm{geo}}$ Closure}

\subsection{The Geometric Prerequisites}
The complete $\Delta_{\mathrm{geo}}$ gap is closed: all five geometric
prerequisites (A9.1--A9.5) are derived from the axioms, not assumed.

\begin{center}
\begin{tabular}{@{}ll@{}}
\toprule
Prerequisite & Derivation \\
\midrule
A9.1 Locality & A1 + finite-bounded cost \\
A9.2 Covariance & T7B \\
A9.3 Conservation & A1 (capacity cannot be created or destroyed) \\
A9.4 Second-order & A4 + Ostrogradsky (no higher-derivative instabilities) \\
A9.5 Propagation & A4 (records require gravitational DOF) \\
\bottomrule
\end{tabular}
\end{center}

\subsection{Lovelock Uniqueness}
\begin{theorem}[Lovelock 1971, as a closure on the APF-prepared geometric class]
In $d=4$, the only symmetric, divergence-free tensor $G_{\mu\nu}$ that depends
on the metric and its first and second derivatives (linearly in second
derivatives) is the Einstein tensor plus a cosmological term.
\end{theorem}

A9.1--A9.5 are precisely Lovelock's conditions. The unique field equation is
\[
G_{\mu\nu} + \Lambda g_{\mu\nu} = \kappa T_{\mu\nu}.
\]
Einstein's equations are not postulated. They are the unique
admissibility-preserving geometric response in $d=4$. The equivalence principle
follows: because correlation cost depends only on capacity distribution (not
internal properties), all systems follow the same least-cost paths. Universality
of free fall follows from universality of correlation constraints. Tag:
[P\_structural] + [I: Lovelock 1971]. Formal closure: \coderef{check\_T9\_grav}{gravity.py}; the unified A9
prerequisite closure is \coderef{check\_A9\_closure}{gravity.py} (v6.9
NEW), which audits all five A9.1--A9.5 prerequisites as a single check
rather than requiring readers to cross-reference \texttt{gravity.py},
\texttt{spacetime.py}, and \texttt{internalization\_geo.py}
individually.

\section{The Cosmological Constant as Capacity Residual (T11)}

\begin{theorem}[T11, cosmological constant as global capacity residual]
$\Lambda$ is the global capacity residual after all local enforcement
commitments:
\[
\Lambda \sim (C_{\mathrm{total}} - C_{\mathrm{used}}) / V.
\]
\end{theorem}
Formal verification: \coderef{check\_T11}{cosmology.py}. The time-independence
claim $w = -1$ at all epochs is proved in \coderef{check\_L\_equation\_of\_state}{cosmology.py};
the DESI DR2 comparison is in \coderef{check\_L\_DESI\_response}{cosmology.py}.

\subsection{Why $\Lambda$ Is Small}
Almost all correlation capacity has been discharged locally: structure
formation, matter, radiation, horizons, records. Only the asymptotic remainder
survives. The remaining scale is set by the largest admissible interface --- the
cosmological horizon $R_H$ --- giving $\Lambda \sim 1/R_H^2 \sim H_0^2$. Not
because of vacuum energy, but because $R_H$ is the last interface where
correlation enforcement still matters.

\subsection{Why $\Lambda$ Is Nonzero}
Exact saturation ($C_{\mathrm{total}} = C_{\mathrm{used}}$) is a measure-zero
condition. Any finite universe with finite total capacity generically has a
small positive residual. The sign is positive because $C_{\mathrm{total}} >
C_{\mathrm{used}}$: the universe has not yet fully committed its capacity
budget.

\subsection{Why $\Lambda$ Looks Like Accelerated Expansion}
Local mass-energy is redistributable: it creates curvature gradients and
attraction. Global residual capacity is non-redistributable: it creates uniform
pressure and expansion. $\Lambda$ does not ``push space apart''; it prevents
further correlation contraction once global slack is exhausted. Acceleration is
the failure mode of further geometric reorganization. Tag: [P\_structural].

\subsection{The APF Capacity Partition Matches Planck Across Four Parameters}\label{sec:planck-match}

The capacity partition of \coderef{check\_L\_saturation\_partition}{cosmology.py}
is not just an internal book-keeping result. Combined with
\coderef{check\_L\_equip}{core.py} (max-entropy distribution), it predicts
the cosmological density parameters of $\Lambda$CDM with zero free parameters:

\begin{center}
\begin{tabular}{@{}lllll@{}}
\toprule
Quantity & APF formula & APF value & Planck 2018 & Difference \\
\midrule
$\Omega_\Lambda$ & $42/61$ & $0.6885$ & $0.6848 \pm 0.007$ & $+0.54\%$ \\
$\Omega_m$       & $19/61$ & $0.3115$ & $0.3152 \pm 0.007$ & $-1.17\%$ \\
$\Omega_b$       & $3/61$  & $0.0492$ & $0.0493 \pm 0.0005$ & $-0.25\%$ \\
$\Omega_{\mathrm{cdm}}$ & $16/61$ & $0.2623$ & $0.2645 \pm 0.006$ & $-0.82\%$ \\
\bottomrule
\end{tabular}
\end{center}

All four parameters match Planck 2018 within $1.2\%$, comfortably inside
Planck's measurement uncertainty. This is a four-parameter prediction
from a single integer partition $61 = 42 + 3 + 16$ that is forced by
gauge-addressability and confinement applied to the anomaly-free
field content (Paper~4). The match is not a fit; the APF integers
are determined entirely by the type-classification predicates and
were not adjusted to cosmological data.

\subsection{Hubble Constant Prediction and the H0DN 2026 Tension}\label{sec:H0-status}
Combining APF's $\Omega_m = 19/61$ with Planck's measured physical matter
density $\Omega_m h^2 = 0.143$ gives a zero-parameter Hubble-constant
prediction:
\[
h_{\mathrm{APF}}^2 = \frac{\Omega_m h^2 \mid_{\mathrm{measured}}}
                          {\Omega_m \mid_{\mathrm{APF}}}
                  = \frac{0.143}{19/61} = 0.4591
\quad\Longrightarrow\quad
H_0^{\mathrm{APF}} = 67.76 \text{ km s}^{-1}\text{Mpc}^{-1}.
\]
This matches the Planck CMB value $67.36 \pm 0.54$ within $1\sigma$.

The H0DN community-consensus measurement (Casertano et al., A\&A 708 A166,
10 April 2026) reports $H_0 = 73.50 \pm 0.81$ km\,s$^{-1}$\,Mpc$^{-1}$
from a covariance-weighted distance-network combination across thirteen
indicator categories. The tension with the APF prediction is
\[
\frac{H_0^{\mathrm{H0DN}} - H_0^{\mathrm{APF}}}
     {\sigma_{\mathrm{H0DN}}}
= \frac{73.50 - 67.76}{0.81} = 7.09\sigma.
\]
DESI DR2 separately favors $w_0 \approx -0.75$, $w_a \approx -0.99$
at $\sim 2\sigma$ from $w = -1$
(\coderef{check\_L\_DESI\_response}{cosmology.py}).

\subsection{Escape Routes and Their Status}\label{sec:H0-routes}
The $H_0$ tension can in principle be resolved by one of five
mechanisms. APF's status against each:

\begin{center}
\begin{tabular}{@{}llll@{}}
\toprule
Route & Mechanism & APF blocker & Empirical viability \\
\midrule
I  & Pre-recombination $r_d$ cut $\sim 7.8\%$ (EDE, extra $N_{\rm eff}$) &
     L\_equation\_of\_state, L\_anomaly\_free &
     EDE is mainstream candidate \\
II & Late-time dynamical DE ($\Omega_\Lambda(a)$ growth $\sim 12\%$) &
     L\_saturation\_partition &
     DESI prefers wrong sign \\
III & Local distance-ladder systematics &
     none &
     H0DN 2026 substantially rules out \\
IV & Spatial curvature $\Omega_k \sim -0.05$ &
     none &
     Planck $|\Omega_k| < 0.02$, insufficient \\
V & Local underdensity $\sim 25\%$ at Hubble-flow scale &
     none &
     Observed $\sim 5$--$15\%$, short by $\sim 2\times$ \\
\bottomrule
\end{tabular}
\end{center}

Routes I and II are blocked by [P] theorems in v6.9. Route III is
substantially ruled out by H0DN itself. Routes IV and V are
APF-admissible (the framework is silent on them) but neither is currently
empirically sufficient. \emph{No route is both APF-admissible and
empirically viable.}

\subsection{Falsifier Statement}\label{sec:H0-falsifier}
The honest disposition under the audit-response discipline:

\begin{apfbox}
\textbf{APF's cosmological-side falsifier.} APF predicts
$H_0 = 67.76 \pm \epsilon$ km\,s$^{-1}$\,Mpc$^{-1}$ (with $\epsilon \sim
0.5$ inheriting Planck's $\Omega_m h^2$ uncertainty), $w = -1$ exactly
at all epochs, and the four-parameter Planck match in
\S\ref{sec:planck-match}. \emph{If the H0DN $H_0 = 73.50 \pm 0.81$
result stands as the true local cosmic expansion rate} --- i.e. is not
explained by Routes III, IV, or V, all of which APF leaves to empirical
determination --- \emph{then at least one of the following [P] theorems
must be demoted in a future codebase version:}
\begin{itemize}[leftmargin=1.5em,nosep]
    \item \coderef{check\_L\_equation\_of\_state}{cosmology.py}
    (would have to admit $w \neq -1$ in some regime),
    \item \coderef{check\_L\_saturation\_partition}{cosmology.py}
    (would have to admit non-topological partition evolution),
    \item \coderef{check\_L\_anomaly\_free}{gauge.py} (would have to
    admit late-time type decoupling).
\end{itemize}
This is named as Project~C.6 in Appendix~\ref{sec:open-A9}.
Phenomenological dynamical-dark-energy proposals that contradict these
checks without first deriving the regime conditions under which they
fail are not adopted into the framework. Any future revision must
discharge the demotion explicitly and prove that the new mechanism
reproduces both the four-parameter Planck match and the H0DN
measurement. The full mathematical analysis is in
\emph{Reference -- H0 Tension Mathematical Analysis (2026-04-18)}.
\end{apfbox}

\section{Status Audit: Dark Matter as $C_{\mathrm{ext}}$}\label{sec:DM}

The capacity partition $C_{\mathrm{total}} = C_{\mathrm{int}} + C_{\mathrm{ext}}$
identifies a geometric sector $C_{\mathrm{ext}}$ that gravitates (it curves
spacetime via T7B) but does not couple to gauge interactions (no internal
structure). This sector naturally invites identification with the astrophysical
dark-matter sector. The identification is attractive and economical, but it
mixes three distinct layers of claim that should be visually fenced apart.

\subsection*{Layer separation}
\begin{enumerate}[leftmargin=1.6em]
    \item \textbf{Structural result.} If the APF partition into internal and
    external capacity sectors is accepted, then the external sector is naturally
    associated with geometric response rather than gauge-charged internal
    excitations. This is a structural reading of the partition itself.
    \item \textbf{Phenomenological identification.} To identify that external
    sector with the astrophysical dark-matter sector requires an additional
    bridge: one must show that the effective stress-energy carried by
    $C_{\mathrm{ext}}$ reproduces the relevant gravitational phenomenology
    across the appropriate scales (rotation curves, lensing, structure
    formation).
    \item \textbf{Empirical discriminator.} The further claim that direct-detection
    experiments should find no dark-matter particles is a downstream empirical
    discriminator, not the same statement as the structural partition.
\end{enumerate}

\begin{statusbox}
\textbf{Recommended status tags.}
\begin{center}
\begin{tabular}{@{}ll@{}}
\toprule
Claim & Recommended tag \\
\midrule
``There is an external geometric capacity sector $C_{\mathrm{ext}}$.'' & [P\_structural] \\
``$C_{\mathrm{ext}}$ provides the correct dark-matter phenomenology.'' & [P\_structural + Pred] \\
``Direct detection should find no dark-matter particles.'' & [Pred] \\
\bottomrule
\end{tabular}
\end{center}
The metric/variational results in this paper are on relatively solid ground as
representation and closure claims, modulo clearly named imports. The dark-matter
identification is more exposed because it crosses from internal structure to
cosmological phenomenology. If presented at the same level as the metric or
Euler--Lagrange results, a hostile reader will attack the whole paper through
its most weakly fenced claim.
\end{statusbox}

\subsection*{Best presentation}
The strongest honest presentation is: the APF capacity partition yields an
external geometric sector that naturally invites identification with
dark-matter-like gravitational load. This motivates a concrete phenomenological
proposal --- that some or all observed dark-matter effects reflect the
$C_{\mathrm{ext}}$ sector rather than a new particle species. The decisive
status of this identification depends on confrontation with lensing,
rotation-curve, and structure-formation data.

\subsection*{Open theorem project}
To upgrade the dark-matter claim, the missing work is not additional rhetoric
but a formal bridge theorem: \emph{given the APF external sector
$C_{\mathrm{ext}}$, derive the effective source term it contributes in the
geometric response equations and prove that its admissible distributions
reproduce at least one nontrivial dark-matter data class without introducing
particle degrees of freedom.} Until such a theorem is in hand, the claim should
remain visibly tagged as a structural proposal with empirical consequences,
not as a closed theorem-level derivation.

\section{Where Least-Cost Geometric Representation Fails}\label{sec:geo-fails}

Geometry is valid when correlation cost is locally integrable, capacity
gradients are smooth, and the admissible class supports a locally quadratic
response structure. Geometry fails when one or more of these hypotheses breaks.

\textbf{Singularities} destroy smooth local cost comparison by driving effective
gradients beyond the regime of local approximation (Type IV exit).

\textbf{Horizons} destroy admissible extendibility: continuation remains
meaningful at the level of admissibility, but no longer as a single smoothly
extendible geometric path from within the original domain (Type II / III
depending on whether the issue is non-extendibility of admissible continuation
or transfer to a new admissible description).

\textbf{Planck-scale discreteness} destroys the tangent-space and
local-additivity assumptions required for metric encoding (Type IV).

\textbf{Topology change} alters the type of admissible class itself, so the
original minimization problem is no longer posed on the same representational
domain (Type IV).

These are not pathologies of the underlying framework. They are predictable
exits from geometric representation. Admissibility remains deeper than geometry;
when geometry fails, the accounting layer survives.

\section{Epistemic Tagging Discipline (Paper 6)}\label{sec:tagging}

\begin{apfbox}
\textbf{Project-wide categories.}
\begin{itemize}[leftmargin=1.5em,nosep]
    \item \textbf{[P] Internally derived.}
    \item \textbf{[P]+[I] Derived with import.} APF prepares the admissible
    class; an external mathematical theorem closes the classification.
    \item \textbf{[P\_structural] Structural identification.}
    \item \textbf{[Pred] Phenomenological proposal/prediction.}
\end{itemize}
\end{apfbox}

\begin{center}
\begin{tabular}{@{}llll@{}}
\toprule
Item & Tag & Coderef & Reason \\
\midrule
Regime R variational structure & [P] & \coderef{check\_Regime\_R}{plec.py} & Smooth, locally additive, connected, unsaturated (v6.9 formalization). \\
Euler--Lagrange form & [P] & --- & Coordinate form of least-cost selection. \\
Metric tensor (T7B) & [P\_structural] & \coderef{check\_T7B}{gravity.py} & Quadratic non-factorizing response. \\
Lorentzian signature & [P\_structural]+[I] & \coderef{check\_L\_HKM\_causal\_geometry}{extensions.py} & APF order + HKM/Malament closure. \\
$d = 4$ (T8) & [P\_structural] & \coderef{check\_T8}{spacetime.py} & Graviton DOF + capacity efficiency. \\
Einstein equations & [P\_structural]+[I] & \coderef{check\_T9\_grav}{gravity.py} & A9 prereqs + Lovelock closure. \\
$\kappa \sim 1/C_{\ast}$ & [P\_structural] & \coderef{check\_T10}{gravity.py} & Scaling-only identification. \\
$\Lambda$ as capacity residual (T11) & [P\_structural] & \coderef{check\_T11}{cosmology.py} & Structural identification/scaling. \\
$w = -1$ at all epochs & [P] & \coderef{check\_L\_equation\_of\_state}{cosmology.py} & Pure cosmological constant. \\
DESI DR2 comparison & [P] & \coderef{check\_L\_DESI\_response}{cosmology.py} & Tension at 1.5--2$\sigma$; revision triggers named. \\
Dark matter as $C_{\mathrm{ext}}$ & [P\_structural + Pred] & --- (open: \S\ref{sec:DM}) & Structural proposal + empirical claim. \\
\bottomrule
\end{tabular}
\end{center}

Whenever a paper summary table contains both theorem-level derivations and
interpretive identifications, the main text should not rhetorically elevate all
entries to the same strength. Imported closure theorems should be described as
closure of an APF-prepared admissible class, not as if APF alone proved the
imported theorem.

\section{Summary}\label{sec:summary}

\subsection{The Derivation Chain}
\begin{itemize}[leftmargin=1.5em,nosep]
    \item Regime R supports a well-defined cost functional; PLEC selects the
    realized path; Euler--Lagrange is its coordinate form.
    \item Conservation laws follow from cost-functional symmetries (Noether).
    \item Distance is minimal correlation cost; geodesics are least-cost paths.
    \item T7B identifies the metric as a non-degenerate symmetric bilinear form
    on tangent spaces.
    \item A4 + HKM/Malament fix Lorentzian signature.
    \item T8 fixes $d = 4$.
    \item A9.1--A9.5 prepare Lovelock's conditions; Lovelock closes uniqueness
    of the Einstein tensor in $d=4$.
    \item T11 identifies $\Lambda$ as global capacity residual.
    \item $C_{\mathrm{ext}}$ provides a structurally motivated dark-matter
    candidate; the identification is fenced as proposal + prediction (\S\ref{sec:DM}).
\end{itemize}

\subsection{The Structural Conclusion}
Dynamics and geometry are not primitive laws but least-cost representational
compressions of admissible correlation reallocation in regimes where the
admissible class is smooth, locally additive, and unsaturated. When those
hypotheses hold, realized paths admit variational and geometric encoding. When
they fail, the representational layer exits its domain while the deeper
admissibility structure remains intact. \emph{Physical law is deeper than its
dynamical or geometric representation.} Together with Paper~5 (quantum
structure), this recovers the full structure of fundamental physics from finite
enforceability --- not as postulates, but as the unique bookkeeping available
when capacity constraints permit smooth description.

\appendix

\section*{Appendix A: $\Delta_{\mathrm{geo}}$ Closure Summary}

The $\Delta_{\mathrm{geo}}$ gap encompassed four sub-gaps, all closed:

\begin{center}
\begin{tabular}{@{}ll@{}}
\toprule
Sub-gap & Resolution \\
\midrule
$\Delta_{\mathrm{ordering}}$ & A4 irreversibility $\Rightarrow$ strict partial order. R2/R3 formalized. \\
$\Delta_{\mathrm{fbc}}$ & Lipschitz lemma from A1 + $\epsilon$-granularity. \\
$\Delta_{\mathrm{continuum}}$ & Kolmogorov + chartability bridge. \\
$\Delta_{\mathrm{signature}}$ & A4 $\to$ order + HKM + $\Omega = 1$. \\
\bottomrule
\end{tabular}
\end{center}

\section*{Appendix B: Enforcement Cost vs. Action}

Entropy measures irreversibly committed capacity (accounting of commitment).
Enforcement cost measures capacity required for transitions (cost of
reallocation). They are related but distinct: entropy accumulates; enforcement
cost is path-dependent. In Regime~R, accumulated enforcement cost is encoded as
action $K = \int L\,dt$. The identification with physical action is exact within
the regime and breaks down at saturation. Paper~7 (ACTION) proves that this
accumulated cost has a minimum: the structural origin of Planck's constant.

\section*{Appendix C: Open Theorem Projects}

The v2.0-PLEC revision identifies open theorem projects directly relevant to
this paper.

\subsection*{Project C.1: PLEC existence and uniqueness conditions}
Define $\mathcal{A}_{\Gamma}$ precisely, formulate $E_{\Gamma}$ on it, identify
the coercivity / lower-semicontinuity analogues that guarantee an argmin
exists, and define the equivalence relation $\sim$ under which uniqueness is
stated. Without this, PLEC remains a suggestive variational slogan rather than
a usable theorem engine.

\subsection*{Project C.2: Regime R as a formal check (CLOSED in v6.9)}\label{sec:open-RegimeR}
\emph{Closed 2026-04-18.} The four R1--R4 conditions are now formalized as
\coderef{check\_Regime\_R}{plec.py}, with an executable witness (1D harmonic
kinetic cost on a connected admissible path class) verifying that R1--R4
jointly imply existence of the PLEC selector. The five regime-exit types
of \S\ref{sec:exit-taxonomy} are formalized as
\coderef{check\_Regime\_exit\_Type\_I}{plec.py} through
\coderef{check\_Regime\_exit\_Type\_V}{plec.py}, each with an explicit
witness instantiating the failure mode. Paper~6 Part~I no longer rests
on R1--R4 as prose conditions; the load-bearing hand-waving risk
flagged in the v2.0-PLEC audit is closed.

\subsection*{Project C.3: A9.1--A9.5 unified-closure check (CLOSED in v6.9)}\label{sec:open-A9}
\emph{Closed 2026-04-18.} The unified Lovelock-prerequisite closure is
now \coderef{check\_A9\_closure}{gravity.py}, which audits A9.1
(locality from A1+FBC), A9.2 (covariance from T7B), A9.3 (conservation
from A1+L\_loc), A9.4 (second-order from A4+Ostrogradsky), A9.5
(propagation from A4+T\_graviton) as a single check rather than
requiring readers to cross-reference three modules.

\subsection*{Project C.4: Formal regime-exit theorem}
Tie each failure mode of \S\ref{sec:dyn-fails} and \S\ref{sec:geo-fails} to a
precise breakdown of a specific PLEC hypothesis. Formalize the Type I--V
taxonomy as a set of checks that each classify a given regime exit by which
hypothesis fails.

\subsection*{Project C.5: Dark-matter bridge theorem}
Derive the effective source term contributed by $C_{\mathrm{ext}}$ in the
geometric response equations and prove that its admissible distributions
reproduce at least one nontrivial dark-matter data class without introducing
particle degrees of freedom (\S\ref{sec:DM}).

\subsection*{Project C.6: $H_0$ tension / dark-energy confrontation}
Address the Casertano et al. (2026) H0DN consensus measurement
($H_0 = 73.50 \pm 0.81$, $7.1\sigma$ from early-universe $\Lambda$CDM) and the
DESI DR2 $w_0/w_a$ tension. Two candidate directions within APF
(\S\ref{sec:H0-status}):
\begin{itemize}
    \item \textbf{Direction A:} Derive what irreversibly committed capacity
    (L\_irr) implies for late-time effective $\Omega_\Lambda(z)$, and check
    whether the result matches DESI's preferred direction ($w_0 > -1$) or
    contradicts it ($w_0 < -1$). Any proposed mechanism must be derived, not
    fit; and must either be consistent with or replace
    \coderef{check\_L\_equation\_of\_state}{cosmology.py}.
    \item \textbf{Direction B:} Explore local-inhomogeneity /
    scale-dependent-enforcement mechanisms that leave $w = -1$ intact but
    address the $H_0$ measurement bias.
\end{itemize}
If neither direction yields a derived mechanism consistent with the data, the
framework accepts an open falsification risk and states it as such.

\section*{Appendix D: Changelog}

\begin{itemize}[leftmargin=1.5em,nosep]
\item \textbf{v2.0-PLEC (April 2026).} Reframed entire derivation around PLEC
schema: admissibility fixes the licensable continuation class; PLEC selects the
realized path; failure modes recast as five formally distinct regime-exit types.
Added shared canonical box, wayfinder box, canonical PLEC schema (definitions),
regime-exit taxonomy section, dark-matter status fence, epistemic tagging
discipline section, open theorem projects appendix, codebase coderef anchors
throughout, and \S\ref{sec:H0-status} on the H0DN/DESI status. Reformatted
bridges around HKM/Malament and Lovelock as imports closing an APF-prepared
admissible class. Cleaned formatting throughout (v1.0 had broken section
spacing and inline math artifacts).
\item \textbf{v1.0 (February 2026).} Complete $\Delta_{\mathrm{geo}}$ closure (T7B,
T8, T11, signature). Dark-matter section. Forward references.
\end{itemize}

\end{document}
